\documentclass[11pt]{article}
\synctex=1
\usepackage[T1]{fontenc}
\usepackage[utf8]{inputenc}
\usepackage{amsmath,amssymb}
\usepackage{geometry}
\usepackage{enumitem}
\usepackage[hidelinks]{hyperref}
\geometry{margin=1in}
\setlist{itemsep=0.25em}

\title{Programmatic Modeling of High-Order Observers for Labeled Finite-State Automata}
\author{Zixi Wang}
\date{}

\begin{document}

\maketitle

\tableofcontents
\newpage

\section{Introduction}

% introduce the problem

Consider a finite-state system $G$ observed by two agents. Agent $A_1$ estimates the
current state of $G$ from its observations. Based on its own observations, Agent
$A_2$ infers the possible state estimates held by Agent $A_1$.

% describe this work

This article develops a programmatic pipeline for constructing the order-2 observer,
which computes Agent $A_2$'s inference about Agent $A_1$'s state estimate.

% transition sentence

The following sections introduce the theoretical foundations and constructions used to implement this pipeline.

\section{Regular Language and FAs}


The following definitions summarize the background needed in this project.

\subsection{Alphabets, words, and languages}


An \textbf{alphabet} is a finite, nonempty set of letters, usually denoted by $\Sigma$. A \textbf{word} over $\Sigma$ is a finite sequence of letters from $\Sigma$. More precisely, a word of length $n$ has the form

\[
w=a_1a_2\cdots a_n
\qquad a_i\in\Sigma.
\]

The empty word is denoted by $\epsilon$. For every $n\in\mathbb{N}$,

\[
\Sigma^0=\{\epsilon\},\qquad
\Sigma^n=\{a_1a_2\cdots a_n\mid a_i\in\Sigma,\ 1\leq i\leq n\}
\quad(n\geq 1).
\]

The \textbf{Kleene star}

\[
\Sigma^*=\bigcup_{n\geq 0}\Sigma^n
=\{\epsilon\}\cup\Sigma\cup\Sigma^2\cup\cdots
\]

is the set of all finite words over $\Sigma$. A \textbf{formal language over $\Sigma$} is any set

\[
L\subseteq\Sigma^*.
\]

Given $L_1,L_2\subseteq\Sigma^*$, there are several basic operations:

\[
\begin{aligned}
L_1\cup L_2
  &=\{w\in\Sigma^*\mid w\in L_1\text{ or }w\in L_2\},\\
L_1\cap L_2
  &=\{w\in\Sigma^*\mid w\in L_1\text{ and }w\in L_2\},\\
\overline{L_1}
  &=\Sigma^*\setminus L_1,\\
L_1\setminus L_2
  &=\{w\in L_1\mid w\notin L_2\},\\
L_1L_2
  &=\{uv\mid u\in L_1,\ v\in L_2\}.
\end{aligned}
\]

For a language $L$, define $L^0=\{\epsilon\}$ and $L^{n+1}=L^nL$. Its \textbf{Kleene star} is

\[
L^*=\bigcup_{n\geq 0}L^n.
\]

\subsection{Regular languages}


A \textbf{regular expression over $\Sigma$} is defined recursively. The expressions $\emptyset$, $\epsilon$, and $a$ for every $a\in\Sigma$ are regular expressions. If $r$ and $s$ are regular expressions, then

\[
(r), \qquad r\cup s,\qquad rs,\qquad r^*
\]

are also regular expressions, representing parentheses,  union, concatenation, and Kleene star, respectively.

The language denoted by a regular expression $r$, written $L(r)$, is defined recursively by

\[
\begin{aligned}
L(\emptyset)&=\emptyset, &
L(\epsilon)&=\{\epsilon\}, &
L(a)&=\{a\},\\
L(r \cup s)&=L(r)\cup L(s), &
L(rs)&=L(r)L(s), &
L(r^*)&=L(r)^*.
\end{aligned}
\]

A language $L\subseteq\Sigma^*$ is \textbf{regular} if and only if there exists a regular expression $r$ over $\Sigma$ such that

\[
L=L(r).
\]

Equivalently, $L$ is regular if and only if it is recognized by a finite automaton.

Regular languages are closed under union, intersection, complement, difference, concatenation, and Kleene star. In other words, if $L_1$ and $L_2$ are regular, then

\[
L_1\cup L_2,\quad
L_1\cap L_2,\quad
\overline{L_1},\quad
L_1\setminus L_2,\quad
L_1L_2,\quad
L_1^*
\]

are also regular.

Let a DFA $A=(Q,\Sigma,\delta,q_0,F)$ recognize the regular language $L$, and choose the \textbf{pumping length}

\[
p=|Q|.
\]

For every $w\in L$ with $|w|\geq p$, the run of $A$ while reading the first $p$ letters visits a state more than once. Therefore,

\[
\forall w\in L\text{ with }|w|\geq p,\;
\exists x,y,z\in\Sigma^*
\]

such that

\begin{itemize}
\item $w=xyz$;
\item $|xy|\leq p$;
\item $y\neq \epsilon$;
\item $\forall i\in\mathbb{N}$, $xy^iz\in L$.

\end{itemize}

Thus, the number of states of a DFA recognizing $L$ is a valid pumping length, although it is not necessarily the smallest possible pumping length.

\subsection{Deterministic and nondeterministic finite automata}


A deterministic finite automaton (DFA) is a quintuple

\[
A=(Q,\Sigma,\delta,q_0,F),
\]

where:

\begin{itemize}
\item $Q$ is a finite set of states;
\item $\Sigma$ is the input alphabet;
\item $\delta:Q\times\Sigma\rightarrow Q$ is the transition function;
\item $q_0\in Q$ is the initial state;
\item $F\subseteq Q$ is the set of final states.

\end{itemize}

The transition function extends from letters to words. For $q\in Q$, $w\in\Sigma^*$, and $a\in\Sigma$,

\[
\delta(q,\epsilon)=q,\qquad
\delta(q,wa)=\delta(\delta(q,w),a).
\]

The DFA accepts $w$ exactly when $\delta(q_0,w)\in F$. Its recognized language is

\[
L(A)=\{w\in\Sigma^*\mid\delta(q_0,w)\in F\}.
\]

A nondeterministic finite automaton (NFA) is a quintuple

\[
N=(Q,\Sigma,\delta,Q_0,F),
\]

where:

\begin{itemize}
\item $Q$ is a finite set of states;
\item $\Sigma$ is the input alphabet;
\item $\delta:Q\times\Sigma\rightarrow 2^Q$ is the transition function;
\item $Q_0\subseteq Q$ is the set of initial states;
\item $F\subseteq Q$ is the set of final states.

\end{itemize}

Here, $2^Q=\{X\mid X\subseteq Q\}$ is the power set of $Q$, so $\delta(q,a)$ may contain zero, one, or several successor states. Extend $\delta$ to sets of states and words by

\[
\delta(X,\epsilon)=X,\qquad
\delta(X,wa)=\bigcup_{q\in\delta(X,w)}\delta(q,a),
\]

for $X\subseteq Q$, $w\in\Sigma^*$, and $a\in\Sigma$. The NFA accepts a word $w$ if at least one state reachable from an initial state is accepting:

\[
\delta(Q_0,w)\cap F\neq\emptyset.
\]

Therefore, the language recognized by $N$ is

\[
L(N)=\{w\in\Sigma^*\mid\delta(Q_0,w)\cap F\neq\emptyset\}.
\]

Every NFA can be converted into an equivalent DFA by the \textbf{subset construction}. 

Kleene's theorem establishes the equivalence

\[
\text{regular expressions}
\quad\Longleftrightarrow\quad
\text{NFAs}
\quad\Longleftrightarrow\quad
\text{DFAs}.
\]

In this project, the underlying system is usually written as

\[
G=(Q,E,\delta,Q_0),
\]

\noindent
where $E$ is an event alphabet, $\delta\subseteq Q\times E\times Q$ is a transition relation, and $Q_0\subseteq Q$ is a set of initial states. Accepting states are omitted because the project studies generated runs and state estimates rather than language acceptance. This transition-system form may be nondeterministic and may have multiple initial states. Adding an observation alphabet and a labeling function to $G$ produces the LFSA introduced in the next section.

\section{LFSAs}


According to Definition 1 in Zhang (2023), a \textbf{labeled finite-state automaton (LFSA)} is a sextuple

\[
\mathcal{S}=(Q,E,\delta,Q_0,\Sigma,\ell),
\]

where:

\begin{itemize}
\item $Q$ is a finite set of states;
\item $E$ is a finite alphabet of events;
\item $\delta:Q\times E\rightarrow 2^Q$ is the transition function;
\item $Q_0\subseteq Q$ is the set of initial states;
\item $\Sigma$ is a finite alphabet of observable labels;
\item $\ell:E\rightarrow\Sigma\cup\{\epsilon\}$ is the labeling function.

\end{itemize}

A transition $q'\in\delta(q,e)$, equivalently $(q,e,q')\in\delta$, means that the occurrence of event $e$ moves the LFSA from state $q$ to state $q'$. Since $\delta(q,e)$ is a set of states, an LFSA may be nondeterministic. 

The event alphabet is partitioned into the observable and unobservable event sets

\[
E_o=\{e\in E\mid\ell(e)\in\Sigma\},
\qquad
E_{uo}=\{e\in E\mid\ell(e)=\epsilon\},
\]

When an event in $E_o$ occurs, its label is observed. When an event in $E_{uo}$ occurs, the observer receives no output.

The transition function is extended from events to finite event sequences. For $q\in Q$, $s\in E^*$, and $e\in E$,

\[
\delta(q,\epsilon)=\{q\},
\qquad
\delta(q,se)=\bigcup_{p\in\delta(q,s)}\delta(p,e).
\]

For $X\subseteq Q$, define

\[
\delta(X,s)=\bigcup_{q\in X}\delta(q,s).
\]

The labeling function is extended from events to event sequences by: 

\[
\ell(\epsilon)=\epsilon,
\qquad
\ell(e_1e_2\cdots e_n)
=\ell(e_1)\ell(e_2)\cdots\ell(e_n).
\]

The set of finite event traces generated by $\mathcal S$ is

\[
L(\mathcal S)
=\{s\in E^*\mid\delta(Q_0,s)\neq\emptyset\}.
\]

The observable language generated by $\mathcal S$ is

\[
\mathcal L(\mathcal S)
=\{\ell(s)\in\Sigma^*\mid s\in L(\mathcal S)\}.
\]

For an observed word $\sigma\in\Sigma^*$, the \textbf{current-state estimate} is

\[
\mathcal M(\mathcal S,\sigma)
=\{q\in Q\mid
\exists q_0\in Q_0,\ 
\exists s\in E^*:
\ell(s)=\sigma
\text{ and }
q\in\delta(q_0,s)\}.
\]

Thus, an LFSA distinguishes an internal event sequence $s$ from its observable label sequence $\ell(s)$. Different event sequences may generate the same observed word, which is the source of state uncertainty studied by concurrent composition and observers.

% src\lfsa.py

% show the python script model for the math model lfsa

\section{Concurrent Composition}

\subsection{Definition}

According to Definition 2 in Zhang (2023), consider two LFSAs with the same output alphabet $\Sigma$:

\[
\mathcal S_i=(Q_i,E_i,\delta_i,Q_{0i},\Sigma,\ell_i),
\qquad i\in\{1,2\}.
\]

For each $\mathcal S_i$, define

\[
E_{io}=\{e\in E_i\mid\ell_i(e)\in\Sigma\},
\qquad
E_{iuo}=\{e\in E_i\mid\ell_i(e)=\epsilon\}.
\]

Their \textbf{concurrent composition} is the LFSA

\[
\operatorname{CC}_{A}(\mathcal S_1,\mathcal S_2)
=(Q',E',\delta',Q'_0,\Sigma,\ell'),
\]

whose components are defined as follows.

\begin{enumerate}
\item The state set is the Cartesian product

\[
   Q'=Q_1\times Q_2.
\]

Thus, a state $(q_1,q_2)$ records the states currently occupied by the two component LFSAs.

\item The initial-state set is

\[
   Q'_0=Q_{01}\times Q_{02}.
\]

\item The observable events synchronize events having the same label:

\[
   E'_o
   =\{(e_1,e_2)\in E_{1o}\times E_{2o}
   \mid\ell_1(e_1)=\ell_2(e_2)\}.
\]

\item The unobservable events occur independently in either component:

\[
   E'_{uo}
   =\{(e_1,\epsilon)\mid e_1\in E_{1uo}\}
   \cup
   \{(\epsilon,e_2)\mid e_2\in E_{2uo}\}.
\]

Therefore,

\[
   E'=E'_o\cup E'_{uo}.
\]

\item For a synchronized observable event $(e_1,e_2)\in E'_o$,

\[
   ((q_1,q_2),(e_1,e_2),(q'_1,q'_2))\in\delta'
   \quad\Longleftrightarrow\quad
   (q_1,e_1,q'_1)\in\delta_1
   \text{ and }
   (q_2,e_2,q'_2)\in\delta_2.
\]

\item For a left-component unobservable event $(e_1,\epsilon)\in E'_{uo}$,

\[
   ((q_1,q_2),(e_1,\epsilon),(q'_1,q'_2))\in\delta'
   \quad\Longleftrightarrow\quad
   (q_1,e_1,q'_1)\in\delta_1
   \text{ and }
   q_2=q'_2.
\]

\item For a right-component unobservable event $(\epsilon,e_2)\in E'_{uo}$,

\[
   ((q_1,q_2),(\epsilon,e_2),(q'_1,q'_2))\in\delta'
   \quad\Longleftrightarrow\quad
   q_1=q'_1
   \text{ and }
   (q_2,e_2,q'_2)\in\delta_2.
\]

\item The labeling function is

\[
   \ell'((e_1,e_2))
   =\ell_1(e_1)
   =\ell_2(e_2),
\]

where the labeling functions are extended by $\ell_1(\epsilon)=\ell_2(\epsilon)=\epsilon$.

\end{enumerate}

If $\mathcal S_1=\mathcal S_2=\mathcal S$, then

\[
\operatorname{CC}_{A}(\mathcal S,\mathcal S)
=\operatorname{CC}_{A}(\mathcal S)
\]

is called the \textbf{self-composition} of $\mathcal S$.

\subsection{Physical Meaning}

Consequently, observable transitions can occur together only when they produce the same label, whereas unobservable transitions can occur separately without changing the other component's state. The concurrent composition therefore collects pairs of runs from $\mathcal S_1$ and $\mathcal S_2$ that generate the same observable label sequence.

\subsection{Applications}

This interpretation makes concurrent composition a general verification tool for discrete-event systems. By pairing runs that are indistinguishable under observation, it reduces questions about detectability, diagnosability, opacity, and related state/event inference or concealment properties to checking property-specific structural conditions on the reachable part of the composed system.

\paragraph{Example---strong detectability.}
After sufficiently many observations, is the current state uniquely known? Self-composition tracks pairs of indistinguishable runs. If a reachable cycle contains a state $(q_1,q_2)$ with $q_1\neq q_2$, ambiguity can persist indefinitely; therefore, strong detectability fails.

\paragraph{Example---fault diagnosability.}
Can a fault always be confirmed after finitely many observations? Concurrent composition pairs a faulty run with a fault-free run that produces the same observation. If a reachable cycle maintains such a pair, the ambiguity can persist indefinitely; therefore, the fault is not diagnosable.

\paragraph{Example---current-state opacity.}
Can an observer determine that the current state is secret? Suppose a secret state $q_s$ and a non-secret state $q_n$ can both generate the same observation $a$. Concurrent composition tracks the indistinguishable pair $(q_s,q_n)$. Because the observation is consistent with a non-secret state, the secret remains opaque. If no such non-secret counterpart is reachable, current-state opacity fails.

\paragraph{Example---event inference.}
Can an observer determine whether a particular event has occurred? Concurrent composition pairs a run containing the event with a run not containing it but producing the same observation. If such a pair is reachable, the event cannot be inferred uniquely; otherwise, its occurrence is identifiable.

\paragraph{Example---secret-event opacity.}
Can the occurrence of a secret event be concealed? Concurrent composition pairs a run containing the secret event with a run containing no secret event but producing the same observation. If every secret run has such an indistinguishable non-secret counterpart, the secret event remains opaque; if a secret run has none, opacity fails.

\section{Observer}

\subsection{Definition}

According to Definition 3 in Zhang (2023), consider an LFSA

\[
\mathcal S=(Q,E,\delta,Q_0,\Sigma,\ell)
\]

with observable event set $E_o$ and unobservable event set $E_{uo}$. Its \textbf{observer} is the deterministic finite automaton

\[
\mathcal S_{\mathrm{obs}}
=(Q_{\mathrm{obs}},\Sigma_{\mathrm{obs}},
\delta_{\mathrm{obs}},q_{0\ \mathrm{obs}}).
\]

Each observer state is a set of possible states of the original LFSA. The components of the observer are defined as follows.

\begin{enumerate}
\item The observer state set is the power set

\[
   Q_{\mathrm{obs}}=2^Q.
\]

\item The observer alphabet is the set of observable labels generated by events:

\[
   \Sigma_{\mathrm{obs}}
   =\ell(E_o)
   =\ell(E)\setminus\{\epsilon\}.
\]

\item For any $X\subseteq Q$, its \textbf{unobservable reach} is

\[
   \operatorname{UR}(X)
   =\bigcup_{q\in X}\ \bigcup_{s\in E_{uo}^*}\delta(q,s).
\]

\item The initial observer state contains every state reachable from an original initial state through only unobservable events:

\[
   q_{0\ \mathrm{obs}}
   =\operatorname{UR}(Q_0).
\]

\item For $X\in Q_{\mathrm{obs}}$ and $a\in\Sigma_{\mathrm{obs}}$, first take every observable event labeled $a$, and then take the unobservable reach:

\[
   \delta_{\mathrm{obs}}(X,a)
   =
   \operatorname{UR}\left(
   \bigcup_{q\in X}
   \ \bigcup_{\substack{e\in E_o\\\ell(e)=a}}
   \delta(q,e)
   \right).
\]

\end{enumerate}

Because $\delta_{\mathrm{obs}}(X,a)$ is one uniquely determined subset of $Q$, the observer is deterministic even when the original LFSA is nondeterministic. For every observed word $\sigma\in\Sigma_{\mathrm{obs}}^*$,

\[
\delta_{\mathrm{obs}}(q_{0\ \mathrm{obs}},\sigma)
=\mathcal M(\mathcal S,\sigma).
\]

\subsection{Physical Meaning}

Thus, after reading $\sigma$, the observer state is exactly the current-state estimate of the LFSA: it contains all and only the states in which $\mathcal S$ could be after producing $\sigma$. Although $Q_{\mathrm{obs}}=2^Q$ contains all subsets theoretically, an implementation usually constructs only the subsets reachable from $q_{0\ \mathrm{obs}}$.

\subsection{Applications}

Observer applications include current-state estimation, detectability verification, fault diagnosis, opacity analysis, and state-based supervisory control under partial observation.

\paragraph{Example---current-state estimation.}
After observing a word $\sigma$, which states may the system currently occupy? Starting from its initial state, the observer follows the transitions labeled by $\sigma$. The reached observer state $\mathcal M(\mathcal S,\sigma)$ contains exactly all possible current states of the original system.

\paragraph{Example---detectability verification.}
Can the current state eventually be determined uniquely from observations? The observer verifies this by examining its reachable state estimates. Persistent reachable estimates containing multiple states reveal observational ambiguity and, depending on the chosen detectability notion, may violate detectability.

\paragraph{Example---fault diagnosis.}
Can the occurrence of an unobservable fault be determined after finitely many observations? Using states augmented with fault information, the observer tracks fault-free and faulty possibilities. An estimate containing both possibilities is uncertain; persistent uncertainty means that the fault is not diagnosable.

\begin{quote}
\textit{Note on fault diagnosis.}
\begin{enumerate}
\item \textit{What is a fault?} A fault is an event, usually unobservable, representing equipment or component failure.
\item \textit{What is a fault-augmented system?} It combines the plant with a memory that records whether a fault has occurred.
\begin{itemize}
\item \textit{What mathematical structure is used to describe the memory?} A two-state automaton $\{N,F\}$, where $N\xrightarrow{f}F$ and $F$ remains $F$.
\end{itemize}
\item \textit{What are plant states and fault statuses?} The plant state $q$ is the current system state, while $N$ or $F$ records whether a fault has occurred. Thus, an augmented state is $(q,N)$ or $(q,F)$.
\end{enumerate}
\end{quote}

\paragraph{Example---opacity analysis.}
Can an external observer determine that the system is in a secret state? Let $Q_S\subseteq Q$ be the secret-state set. If a reachable observer state $X$ satisfies $X\subseteq Q_S$, the observation reveals the secret; therefore, current-state opacity fails.

\paragraph{Example---state-based supervisory control.}
Which controllable events can be enabled safely under partial observation? For each observer state, a supervisor considers every possible current state represented by that estimate and enables only decisions that satisfy the control requirement for all of them.

\begin{quote}
\textit{Note on state-based supervisory control.}
\begin{itemize}
\item \textit{What is a controllable event?} The event set is partitioned as $E=E_c\mathbin{\dot\cup}E_{uc}$. Events in $E_c$ may be disabled by the supervisor, whereas events in $E_{uc}$ cannot.
\item \textit{What is a forbidden state?} Forbidden states form a predefined set $Q_{\mathrm{bad}}\subseteq Q$ containing states that the controlled system must avoid.
\end{itemize}
\end{quote}

\section{High-order observer}

\subsection{Model of Agent}

Consider an FSA

\[
G=(Q,E,\delta,Q_0).
\]

An \textbf{agent} $A_i$ is modeled by the triple

\[
A_i=(E_i,\Sigma_i,\ell_i),
\]

where:

\begin{itemize}
\item $E_i\subseteq E$ is the set of events observable to $A_i$;
\item $\Sigma_i$ is the finite alphabet of labels observed by $A_i$;
\item $\ell_i:E_i\rightarrow\Sigma_i$ is the labeling function of $A_i$.
\end{itemize}

The labeling function is extended to $\bar{\ell}_i:E\rightarrow\Sigma_i\cup\{\epsilon\}$ by

\[
\bar{\ell}_i(e)=
\begin{cases}
\ell_i(e), & e\in E_i,\\
\epsilon, & e\in E\setminus E_i.
\end{cases}
\]

For simplicity, the extended function $\bar{\ell}_i$ is also denoted by $\ell_i$ in the following. Thus, an event in $E_i$ produces an observable label, while an event outside $E_i$ is unobservable to $A_i$. Different observable events may have the same label.

For an event trace $s=e_1e_2\cdots e_n\in E^*$, the observation made by $A_i$ is

\[
\ell_i(s)=\ell_i(e_1)\ell_i(e_2)\cdots\ell_i(e_n)\in\Sigma_i^*,
\]

where every occurrence of $\epsilon$ disappears under concatenation. Two event traces are indistinguishable to $A_i$ if they produce the same observation:

\[
s\sim_{A_i}t
\quad\Longleftrightarrow\quad
\ell_i(s)=\ell_i(t).
\]

Combining the FSA $G$ with the observation mechanism of $A_i$ produces the LFSA

\[
G_{A_i}=(Q,E,\delta,Q_0,\Sigma_i,\ell_i).
\]

The agent does not change the states, events, or transitions of $G$; it only determines how the events are observed. Consequently, different agents may produce different LFSAs from the same FSA because they may have different observable event sets, alphabets, or labeling functions.

\subsection{Definition}

Consider an FSA

\[
G=(Q,E,\delta,Q_0)
\]

\noindent
and an ordered sequence of agents $A_1,A_2,\ldots,A_n$. Agent $A_i$ observes an event trace $s\in L(G)$ through its labeling function

\[
\ell_i:E\rightarrow\Sigma_i\cup\{\epsilon\}.
\]

An ordinary observer gives an agent's estimate of the current state of $G$. This is an \textbf{order-1 estimate}:

\[
\mathcal M_{A_1}(G,\alpha)
=\{q\in Q\mid
\exists q_0\in Q_0,\ 
\exists s\in L(G):
\ell_1(s)=\alpha
\text{ and }
q\in\delta(q_0,s)\}
\in 2^Q.
\]

A high-order estimate describes one agent's estimate of another agent's estimate. For example, after $A_2$ observes $\beta\in\ell_2(L(G))$, its order-2 estimate of what $A_1$ may know is

\[
\mathcal M_{A_1\leftarrow A_2}(G,\beta)
=
\{\mathcal M_{A_1}(G,\ell_1(s))
\mid
s\in L(G),\ \ell_2(s)=\beta\}
\in 2^{2^Q}.
\]

\subsection{Physical Meaning}

Thus, $A_1$ estimates the current state of $G$, whereas $A_2$ estimates the possible state estimates held by $A_1$. The arrow in $A_1\leftarrow A_2$ indicates that the final inference is made by $A_2$ about $A_1$.

\subsection{Applications}

High-order observers provide a systematic way to solve \textbf{epistemic inference problems} in discrete-event systems. They can determine what one agent knows about another agent's state estimate, verify higher-order opacity and information-flow privacy, analyze epistemic detectability and diagnosability, and support decentralized supervisory-control decisions that depend on nested knowledge.

\paragraph{Example---high-order opacity.}
Can agent $A_2$ determine that agent $A_1$ knows the system is in a secret state? The order-2 observer tracks all state estimates that $A_1$ may hold from $A_2$'s viewpoint. If a reachable order-2 estimate contains only $A_1$-estimates that lie entirely within the secret-state set, then $A_2$ knows that $A_1$ knows the secret; therefore, high-order opacity fails.

\paragraph{Example---information-flow privacy.}
Can agent $A_2$ infer private information available to agent $A_1$? The high-order observer tracks the estimates that $A_1$ may hold after each observation of $A_2$. If a reachable high-order estimate reveals only one possible private value, then $A_2$ can infer that value; therefore, information-flow privacy fails.

\begin{quote}
\textit{Note on information-flow privacy.}
\begin{itemize}
\item \textit{What is private information?} It is an attribute of $A_1$'s state estimate that should not be uniquely inferred by $A_2$, such as whether the estimate reveals a secret state, event, or control decision.
\item \textit{How is private information represented?} Let $h:2^Q\rightarrow V$, where $V$ is the set of possible private values; for example, $V=\{\text{secret},\text{non-secret}\}$.
\item \textit{What does $A_1$ hide?} Agent $A_1$ hides the private attribute value $h(X)$ associated with its state estimate $X$, rather than necessarily hiding the state estimate itself.
\end{itemize}
\end{quote}

\paragraph{Example---epistemic detectability.}
Can agent $A_2$ determine that agent $A_1$ knows the current state uniquely? The order-2 observer tracks all estimates that $A_1$ may hold from $A_2$'s viewpoint. If every estimate in a reachable order-2 state is the same singleton, then $A_2$ knows that $A_1$ knows the current state.

\begin{quote}
\textit{Note on the same singleton.}
If a reachable order-2 state is $\{X_1,\ldots,X_k\}$, then ``every estimate is the same singleton'' means
\[
X_1=\cdots=X_k=\{q\}
\]
for some $q\in Q$. Thus, all possible estimates held by $A_1$ agree on exactly one current state.
\end{quote}

\paragraph{Example---epistemic diagnosability.}
Can agent $A_2$ determine that agent $A_1$ has diagnosed a fault? First augment the plant states with a fault status. The order-2 observer then tracks all augmented-state estimates that $A_1$ may hold. If every estimate in a reachable order-2 state contains only states marked as faulty, then $A_2$ knows that $A_1$ has diagnosed the fault.

\paragraph{Example---decentralized supervisory control.}
Can agent $A_2$ safely enable an event whose decision depends on what agent $A_1$ knows? The high-order observer collects all estimates that $A_1$ may hold from $A_2$'s viewpoint. Agent $A_2$ enables the event only if the required control condition holds for every such estimate; otherwise, enabling it may violate the specification.

\subsection{Pipeline to get the high-order observer}


Consider an FSA $G=(Q,E,\delta,Q_0)$ and two ordered agents
$A_1=(E_1,\Sigma_1,\ell_1)$ and
$A_2=(E_2,\Sigma_2,\ell_2)$. An event outside $E_i$ is
unobservable to $A_i$, so its label under $\ell_i$ is
$\epsilon$. The order-2 observer is constructed as follows.

\begin{enumerate}
\item Combine $G$ with $A_1$'s observation mechanism to obtain the
LFSA

\[
   G_{A_1}=(Q,E,\delta,Q_0,\Sigma_1,\ell_1).
\]

In the program, \texttt{src/fa\_agent\_to\_lfsa.py} reads the FSA JSON and the
$A_1$ JSON and writes the LFSA JSON.

\item Apply the powerset construction to $G_{A_1}$:

\[
   \operatorname{Obs}_{\ell_1}(G)
   =:\operatorname{Obs}_{A_1}.
\]

Each observer state $X\subseteq Q$ is $A_1$'s current-state
estimate after an observation. The program \texttt{src/observer.py}
constructs the nonempty reachable part of this observer. Although
the observer is mathematically a DFA, its JSON representation uses
identity labels, $\ell_{\mathrm{obs}}(a)=a$, so it can be passed
to the concurrent-composition program.

\item Concurrently compose $G_{A_1}$ and
$\operatorname{Obs}_{A_1}$:

\[
   CC_{\ell_1}
   \left(G_{A_1},\operatorname{Obs}_{A_1}\right).
\]

A state of this composition has the form $(q,X)$, where $q$ is
the actual state of $G$ and $X$ is the state estimate held by
$A_1$. Observable transitions synchronize under the same
$\ell_1$-label, while unobservable transitions of $G_{A_1}$
interleave. This step is implemented by \texttt{src/CC.py}.

\item The composition is still labeled from $A_1$'s perspective. To
prepare it for $A_2$, first project every paired composition event

\[
   (e_1,e_2)\longmapsto e_1
\]

and remove the existing $\ell_1$-labels. This produces the
underlying FSA whose states remain pairs $(q,X)$. The operation is
provided by \texttt{src/lfsa\_to\_fa.py} using its \texttt{--project-left} option.

\item Combine this projected FSA with $A_2$. This replaces
$\ell_1$ with $\ell_2$ and produces the LFSA
$CC^{G,\mathrm{Obs}}_{A_1\rightarrow A_2}$. In the program, run
\texttt{src/fa\_agent\_to\_lfsa.py} with the projected FSA and the $A_2$
agent JSON.

\item Apply \texttt{src/observer.py} to the $A_2$-labeled composition:

\[
   \operatorname{Obs}_{\ell_2}
   \left(CC^{G,\mathrm{Obs}}_{A_1\rightarrow A_2}\right)
   =:
   \operatorname{Obs}_{A_1\leftarrow A_2}(G).
\]

This is the order-2 observer. A state has the form

\[
   \{(q_1,X_1),\ldots,(q_n,X_n)\}.
\]

After $A_2$ observes a word $\alpha$, its possible inferences of
$A_1$'s estimate are

\[
   M_{A_1\leftarrow A_2}(G,\alpha)
   =\{X_1,\ldots,X_n\}.
\]

\end{enumerate}

The programmatic data flow is therefore

\begin{verbatim}
FSA G + Agent A1
        |
        v
LFSA G_A1
        |
        v observer.py
Obs_A1 (DFA stored with identity labels)
        |
        +-------------------+
        |                   |
        v                   |
CC.py <----- LFSA G_A1 -----+
        |
        v
CC labeled by A1
        |
        v lfsa_to_fa.py --project-left
Projected underlying FSA
        |
        + Agent A2
        v fa_agent_to_lfsa.py
CC labeled by A2
        |
        v observer.py
Order-2 observer Obs_A1<-A2(G)
\end{verbatim}


For $n>2$, the construction is repeated recursively: project the
current composition onto the original event component, apply the next
agent's labeling function, and construct another observer. Each new
observer adds one level of state-estimate nesting. In the general case,
constructing an order-$n$ observer has $n$-EXPTIME complexity.

\section{Future Work}

Future work will focus on the following directions:

\begin{itemize}
\item automate the recursive construction of order-$n$ observers;
\item add verification procedures for high-order opacity, detectability, and diagnosability;
\item develop benchmark examples and automated correctness tests;
\end{itemize}


\addcontentsline{toc}{section}{References}
\begin{thebibliography}{9}

\bibitem{zhang2023concurrent}
K.~Zhang,
``A unified concurrent-composition method to state/event inference and concealment in labeled finite-state automata as discrete-event systems,''
\emph{Annual Reviews in Control},
vol.~56, article~100902, 2023.
\url{https://doi.org/10.1016/j.arcontrol.2023.100902}.

\bibitem{zhang2024highorder}
K.~Zhang, X.~Han, A.~Giua, and C.~Seatzu,
``[Draft] High-order observers and high-order state-estimation-based properties of discrete-event systems,''
\emph{arXiv preprint} arXiv:2408.06141, 2024.
\url{https://doi.org/10.48550/arXiv.2408.06141}.

\end{thebibliography}

\end{document}
